% ─────────────────────────────────────────────────────────────────────────────
% METHODOLOGY SECTION --- clean rewrite
% Principles:
%   - Operational focus: equations + implementation choices, not re-introduction
%   - Related Work already introduced EWC, OGD, FOPNG, HN conceptually
%   - Novel contributions (iFOPNG, normalization, dropout decision) get full
%     treatment; baselines get equations + key parameters only
%   - ONG and FNG filled in from Garg et al. (2026) formulations
%   - Dropout paragraph cleaned of HTML rendering artifacts
% ─────────────────────────────────────────────────────────────────────────────

\section{Methods}
\label{sec:methods}
% ─────────────────────────────────────────────────────────────────────────────
\subsection{Continual Learning Protocol}
\label{sec:cl_protocol}

\subsubsection{Task Setup and Evaluation Metrics}
A continual learning agent receives a sequence of $T$ tasks
$\mathcal{T}_1, \mathcal{T}_2, \ldots, \mathcal{T}_T$ presented in order.
At task $t$, only the data from $\mathcal{T}_t$ is available; prior task
data is not retained.
After each task is trained, the model is evaluated on all tasks seen so far.
We report two aggregate metrics: \emph{average accuracy}
\begin{equation}
    A_t = \frac{1}{t}\sum_{i=1}^{t} a_{t,i},
\end{equation}
and \emph{backward transfer}
\begin{equation}
    \mathrm{BWT} = \frac{1}{T-1}\sum_{t=2}^{T}(a_{T,t}-a_{t,t}),
\end{equation}
where $a_{i,j}$ denotes accuracy on task $j$ immediately after training
task $i$.
Negative BWT quantifies catastrophic forgetting; a value of zero indicates
perfect retention.

Reporting both metrics jointly is necessary because each can mask
failure modes invisible to the other: a method that forecloses new
learning to preserve old tasks achieves near-zero BWT alongside collapsed
accuracy, while a method that overfits each new task sustains competitive
accuracy at the cost of large negative BWT on earlier ones.

% ─────────────────────────────────────────────────────────────────────────────
\subsection{Fisher Information Matrix and Regularization-Based Continual Learning}
\label{sec:fisher_ewc}

All Fisher-based methods evaluated in this thesis use the diagonal empirical
approximation $F$ defined in Equation~\eqref{eq:empirical_fisher}.
The $F_{\mathrm{new}}$ and $F_{\mathrm{old}}$ used in
projection methods denote the FIMs computed from the current task's data and
accumulated from all prior tasks, respectively; both are instances of Equation~\eqref{eq:empirical_fisher}.

\subsubsection{Elastic Weight Consolidation}
\label{sec:ewc}
EWC (Section~\ref{sec:rw_regularization}, Equation~\ref{eq:ewc}) is included as
a regularization baseline that operates purely through parameter-space
penalties and maintains no gradient subspace.
We use $\lambda = 10^{-3}$ for all hypernetwork configurations and
reproduce the per-benchmark values from \citet{garg_fisher-orthogonal_2026}
Table~1 for standalone evaluations.
Fisher estimates are computed over the full training set for MNIST-based
benchmarks and over $N{=}1024$ samples for CIFAR benchmarks.

% ─────────────────────────────────────────────────────────────────────────────
\subsection{Gradient Projection Methods}
\label{sec:projection_methods}

This subsection presents the full suite of gradient-based methods evaluated
in this thesis.
\texttt{SGD} and \texttt{FNG} establish unconstrained lower bounds with no
forgetting-protection mechanism; \texttt{OGD} through \texttt{iFOPNG} form
the projection family that is the primary focus of this work.
% Basic learning algorithm
\subsubsection{Stochastic Gradient Descent}
\texttt{SGD} applies vanilla gradient descent with no forgetting protection:
\begin{equation}
    \Delta\theta_{\mathrm{SGD}} = -\eta\,\nabla_\theta\mathcal{L}_t.
\end{equation}
\texttt{Adam} \citep{kingma_adam_2015} is included as an adaptive-rate
variant that maintains per-parameter momentum; both serve as unconstrained
baselines.

% Basic projection, Euclidean space
\subsubsection{Orthogonal Gradient Descent}
\label{sec:ogd}
OGD \citep{farajtabar_orthogonal_2019} maintains a memory matrix
$G = [g_1 \mid \cdots \mid g_K] \in \mathbb{R}^{p \times K}$ of past task
gradient directions and updates parameters via the current gradient
projected onto the orthogonal complement of $\operatorname{Col}(G)$:
\begin{equation}
    \tilde{g} = g - G(G^\top G)^{-1}G^\top g,
    \qquad
    \Delta\theta_{\mathrm{OGD}} = -\eta\,\tilde{g}.
\end{equation}

\paragraph{Gradient memory management.}
In all projection methods, $G$ is maintained as a fixed-budget basis of at most $M$ columns; 
insertion, orthonormalisation, and overflow recompression are detailed in Section~\ref{sec:normalization}.

% Riemannian step without projection
\subsubsection{Natural Gradient and Fisher Natural Gradient}
Natural gradient descent \citep{amari_natural_1998} replaces the Euclidean
gradient step with one measured in the Riemannian geometry of the Fisher
metric, so that equal step sizes correspond to equal changes in model
predictions rather than equal changes in parameter values:
\begin{equation}
    \Delta\theta_{\mathrm{nat}} = -\eta\,F^{-1}_{\mathrm{new}}\nabla_\theta\mathcal{L}.
\end{equation}
\texttt{FNG} \citep{garg_fisher-orthogonal_2026} applies a trust-region
constraint to this step, normalising by the Fisher norm of the gradient:
\begin{equation}
    v^*_{\mathrm{FNG}}
    = -\eta\,\frac{F^{-1}_{\mathrm{new}} g}
                  {\sqrt{g^\top F^{-1}_{\mathrm{new}} g}}.
\end{equation}
FNG maintains no gradient subspace and provides no forgetting protection;
it is included to isolate the contribution of Fisher preconditioning
independently of projection.

% OGD projection applied to the natural gradient step
\subsubsection{Orthogonal Natural Gradient}
ONG \citep{yadav_ong_2025} extends OGD by replacing the raw gradient with
the natural gradient before projection.
The natural gradient $F^{-1}_{\mathrm{new}}g$ is first computed, then
projected onto the Euclidean orthogonal complement of $\operatorname{Col}(G)$:
\begin{equation}
    \Delta\theta_{\mathrm{ONG}}
    = -\eta\!\left(
        F_{\mathrm{new}}^{-1}g
        - G(G^\top G)^{-1}G^\top F_{\mathrm{new}}^{-1}g
    \right).
\end{equation}
The step direction is Fisher-preconditioned, but the orthogonality
constraint is still measured in Euclidean space: the kernel being
inverted is $A = G^\top G + \lambda I$, identical to OGD.
Following the diagonal Fisher approximation used throughout this thesis,
we substitute the EKFAC inverse of \citet{yadav_ong_2025} with the
diagonal empirical Fisher; results are not directly comparable to the % ITS ALREADY STATED SOMEWHERE
original implementation.

% Both step and projection in Riemannian space
\subsubsection{Fisher-Orthogonal Projected Natural Gradient Descent}
\texttt{FOPNG} \citep{garg_fisher-orthogonal_2026} moves both the step
and the projection into the Riemannian geometry of the Fisher metric.
Where ONG projects the natural gradient in Euclidean space, FOPNG measures
orthogonality in Fisher space: a direction is orthogonal to the stored
basis if it produces no change in model predictions along those directions,
not merely no change in parameter values.
The constrained trust-region update is:
\begin{equation}
    v^*_{\mathrm{FOPNG}} = -\eta\,
    \frac{F^{-1}_{\mathrm{new}} Pg}{\sqrt{(Pg)^\top F^{-1}_{\mathrm{new}} Pg}},
    \qquad
    P = I - F_{\mathrm{old}}\, G\!\left(G^\top F_{\mathrm{old}}\,
        F^{-1}_{\mathrm{new}} F_{\mathrm{old}}\, G
        + \lambda I\right)^{\!-1}\! G^\top F_{\mathrm{old}},
    \label{eq:fopng}
\end{equation}
where $F_{\mathrm{new}}$ is the current-task empirical Fisher,
$F_{\mathrm{old}}$ is the accumulated historical Fisher under which the
stored directions were collected, and $\lambda I$ is a damping term for
numerical stability.
This is the pre-Fisher form of \citet{garg_fisher-orthogonal_2026}, in
which the stored basis carries its collection-time Fisher weighting.

% Single substitution: F_new -> F_c
\subsubsection{Inertial FOPNG}
\label{sec:ifopng}

In FOPNG the Riemannian metric governing both the trust-region constraint
and the projection objective is $F_{\mathrm{new}}$ alone.
This is natural on the first task but creates an asymmetry as history
accumulates: the subspace $\operatorname{Col}(G)$ encodes directions
important to prior tasks, yet the metric assigns zero curvature to
parameters that were important previously but are uninformative on the
current task, allowing the update to move freely through historically
sensitive directions.

\texttt{iFOPNG} resolves this asymmetry by replacing $F_{\mathrm{new}}$
with the \emph{combined metric}
\begin{equation}
    F_c \;\coloneqq\; F_{\mathrm{new}} + F_{\mathrm{old}},
    \label{eq:combined_fisher}
\end{equation}
where $F_{\mathrm{old}}$ is maintained as the element-wise maximum of all
per-task Fisher diagonals accumulated so far.
Substituting $F_c$ for $F_{\mathrm{new}}$ throughout
Equation~\eqref{eq:fopng} yields:
\begin{equation}
    v^* = -\eta\,
    \frac{F_c^{-1}Pg}{\sqrt{(Pg)^\top F_c^{-1}Pg}},
    \qquad
    P = I - F_{\mathrm{old}}\,G\,A^{-1}\,G^\top F_{\mathrm{old}},
    \label{eq:ifopng}
\end{equation}
where $A = G^\top F_{\mathrm{old}}\,F_c^{-1}\,F_{\mathrm{old}}\,G
+ \lambda I$.
This update is derived as the closed-form solution of the trust-region
projection problem in Theorem~\ref{thm:ifopng}
(Appendix~\ref{app:a}).

The combined metric assigns each coordinate $i$ an effective curvature
$(F_{\mathrm{new}})_i + (F_{\mathrm{old}})_i$, suppressing the step size
in directions important to any task in the sequence.
We call this \emph{parameter inertia}: a coordinate that accumulated high
Fisher importance is not pulled back toward a stored reference point ---
it is made harder to displace.
This contrasts with EWC \citep{kirkpatrick_overcoming_2017}, which achieves
retention through an explicit quadratic penalty anchored at $\theta^*$;
\texttt{iFOPNG} achieves it geometrically, through the ambient metric alone.

% ─────────────────────────────────────────────────────────────────────────────
\subsection{Hypernetwork Architecture}
\label{sec:hn_arch}

\subsubsection{Chunked Hypernetworks}
Generating $\theta_t$ in a single forward pass would require an output
layer of dimension $|\theta_t|$, which is computationally prohibitive for
any non-trivial target architecture.
\emph{Weight chunking} \citep{oswald_continual_2020} resolves this by
iterating over a secondary chunk embedding $c_k$, $k=1,\ldots,K$,
and producing a block of $C_s$ target parameters per pass; the full weight
vector is assembled by concatenation.
Chunking makes the hypernetwork parameter-efficient but introduces the
chunk-induced gradient accumulation pathology addressed in
Section~\ref{sec:normalization}.

\subsubsection{Functional Regularizer}
The functional regularizer of \citet{oswald_continual_2020} is active for
all hypernetwork configurations, with regularization strength $\beta$ treated
as a tunable hyperparameter (see Appendix~\ref{app:hyperparameters}).

\subsubsection{Target Network and Dropout}
\label{sec:dropout}
Standalone configurations reproduce the target network architecture of
\citet{garg_fisher-orthogonal_2026} without modification, including two
\texttt{Dropout}$(p{=}0.5)$ layers in the \texttt{SimpleCIFARCNN} classifier
head.
Dropout is removed from the target network in all hypernetwork configurations.
In chunked generation, the hypernetwork executes $K$ sequential forward
passes through the same target architecture before any gradient update;
because all passes share the parent module's training state, stochastic
dropout masks are resampled independently on every chunk pass, injecting
uncontrolled noise into the gradient signal accumulated across chunks.
Removing dropout eliminates this source of variance without modifying
any other aspect of the target architecture
% ─────────────────────────────────────────────────────────────────────────────
\subsection{Novel Algorithmic Contributions}
\label{sec:novel}

\subsubsection{Projection-Scoped Normalization}
\label{sec:normalization}

When a hypernetwork spawns $K$ sequential weight-chunks before a single
projection step executes, the accumulated gradient vector and the diagonal
Fisher entries grow in magnitude proportionally to $K$, inflating the
condition number of the projection kernel $A$ and causing the
matrix inversion in Equation~\eqref{eq:fopng} to fail catastrophically for
any practically sized target network.

We resolve this with two independent rescaling operations that address
distinct sources of ill-conditioning:

\begin{enumerate}
    \item \textbf{Gradient basis orthonormalization.}
        Each batch of $K$ gradient vectors collected after task $t$ is
        replaced by the orthonormal factor $Q$ from its thin QR
        decomposition before insertion into $G$:
        \begin{equation}
            [g_{t,1} \mid \cdots \mid g_{t,K}]
            \;\xrightarrow{\;\mathrm{QR}\;}\; Q,
            \qquad Q^\top Q = I.
        \end{equation}
        QR is chosen over SVD at insertion time because the operation
        requires only orthonormalization of the incoming batch --- the
        number of columns is preserved, not reduced, so the ordered
        factorization that SVD provides is unnecessary overhead.
        SVD is reserved for the overflow handler: when the column count
        of $G$ exceeds the budget $M$, the full basis is recompressed as
        $G \leftarrow U_{:,\,1:M}$, retaining the $M$ left singular
        vectors of greatest variance and discarding the rest.
        This operation applies \emph{at storage time} and affects only
        the basis $G$; the gradient $g$ entering the projection update
        itself is the unmodified parameter gradient from the model.

    \item \textbf{Metric-tensor scaling.}
        Within each projection computation, all Fisher quantities are
        jointly divided by a single scale factor derived from the
        maximum entry of the method's ambient metric $F^*$:
        \begin{equation}
            s = \max\bigl(F^*,\; 1\bigr),
            \qquad
            \tilde{F} = F \,/\, s
            \;\;\text{for } F \in \bigl\{F_{\mathrm{old}},\,
            F_{\mathrm{new}},\, F_c\bigr\},
            \qquad \tilde{\lambda} = \lambda \,/\, s,
            \label{eq:scaling}
        \end{equation}
        where $F^* = F_{\mathrm{new}}$ for \texttt{FNG}, \texttt{ONG},
        and \texttt{FOPNG}, and $F^* = F_c$ for \texttt{iFOPNG}.
        The clamp to 1 prevents division below unity when the Fisher is
        near-zero.
        The damping coefficient $\lambda$ is rescaled by the same factor,
        so its ratio to the Fisher magnitude is held fixed as the absolute
        Fisher scale drifts across tasks. This is necessary because a
        single additive $\lambda$ cannot regularise the kernel by itself
        once chunk accumulation has inflated its spectrum: with the
        smallest and largest eigenvalues of $A$ separated by as much as ten
        orders of magnitude (Section~\ref{sec:results_rq2}), any $\lambda$
        large enough to lift the smallest eigenvalues away from zero
        simultaneously swamps the informative mid-spectrum directions,
        while any $\lambda$ small enough to leave those directions intact
        is negligible against the largest eigenvalues and does nothing for
        the conditioning. Because the condition number is invariant to a
        positive scalar multiple of $A$, the rescaling does not reduce
        $\kappa(A)$ directly; its purpose is to return the absolute
        magnitudes of $A$, its damped inverse, and the trust-region
        denominator to a range in which finite-precision arithmetic
        retains its significant digits and the fixed $\lambda$ acts as a
        meaningful regulariser rather than as numerical noise. 
        % The normalisation is heuristic: it stabilises the inversion of $A$
        % without providing algebraic guarantees on the resulting update % Hmm, maybe thats not needed.
        direction.
\end{enumerate}

Together, these operations ensure that the method-specific projection
kernel $A$ remains well-conditioned regardless of the number of chunks
spawned per task.
The kernel differs structurally across methods:
\begin{equation}
    A =
    \begin{cases}
        G^\top G + \lambda I
            & \text{\texttt{OGD}, \texttt{ONG}} \\[4pt]
        G^\top F_{\mathrm{old}}\,
            F_{\mathrm{new}}^{-1}\,
            F_{\mathrm{old}}\, G + \lambda I
            & \text{\texttt{FOPNG}} \\[4pt]
        G^\top F_{\mathrm{old}}\,
            F_c^{-1}\,
            F_{\mathrm{old}}\, G + \lambda I
            & \text{\texttt{iFOPNG}}
    \end{cases}
    \label{eq:projection_kernel}
\end{equation}
Gradient orthonormalization eliminates column-scale inflation in $G$,
which directly reduces the spectral range of $A$ for all three variants.
Metric-tensor normalization bounds $\tilde{F}$ to $[0,\,1]$, preventing
the Fisher-weighted variants from producing numerically singular $A$.
The \texttt{OGD}/\texttt{ONG} kernel is unaffected by metric-tensor
scaling by construction.
Subspace membership is invariant to positive scalar rescaling of basis
vectors, so neither operation alters the geometry of the projection.

% \paragraph{Principled rescaling and open problem.}
% \label{sec:ng_limitation}
% The proposed normalization --- QR orthonormalization of the gradient basis
% combined with max-entry scaling of $F_c$ --- is heuristic in nature: it
% discards the eigenvalue structure of the pullback metric and applies
% isotropic rather than curvature-aware rescaling to the chunk embedding
% directions.
% The principled alternative is the natural gradient under the Riemannian
% metric induced by the chunk-generating map:
% \begin{equation}
%     \Delta c_i = -\eta\,(J_i^\top J_i + \lambda I)^{-1}\,
%                   \nabla_{c_i}\mathcal{L},
%     \qquad
%     J_i = \frac{\partial W_i}{\partial c_i} \in \mathbb{R}^{C_s \times d_c},
% \end{equation}
% which rescales each chunk embedding direction in proportion to how
% strongly it moves the generated weights, and is tractable because
% $d_c \ll C_s$.
% This preconditioner was implemented and evaluated experimentally.
% Applied to the raw gradient \emph{prior} to the projection step, it
% consistently degraded continual learning performance --- BWT worsened by
% approximately $0.33$ on Split-MNIST SH relative to the heuristic baseline.
% We attribute this to a metric mismatch: the FOPNG projection operator is
% calibrated for gradients in Euclidean parameter space, and preconditioning
% only the chunk embedding component before projection produces a
% geometrically inconsistent input to the kernel $G^\topFG$.
% The principled resolution --- applying $(J_i^\top J_i + \lambda I)^{-1}$
% to the projected update $v^*$ rather than to the raw gradient, so that the
% projection sees Euclidean gradients and the Riemannian step follows --- is
% left to future work.